\documentclass{article}

\usepackage{ctex}
\usepackage{geometry}
\geometry{a4paper, scale=0.85}
\usepackage{graphicx} % 插入图片
\graphicspath{{D:/文章中枢/课程讲义/模拟电子技术-基于Multisim实现/Image}} % 图片路径
\usepackage{subcaption} 
\usepackage{amsmath} % 数学公式
\usepackage{amssymb} % 数学符号
\usepackage{fancyhdr} % 页眉页脚
\usepackage{circuitikz} % 电路图
\circuitikzset{european resistors}
\usepackage{tikz} % 绘图
\usepackage{float} % 图片浮动
\usepackage{listings} % 代码高亮
\usepackage{color} % 颜色
\usepackage{hyperref} % 超链接
\usepackage{siunitx} % SI单位 
\usepackage{comment}
\usepackage{caption} % 图片标题
\captionsetup{labelsep=period, format=hang}
\numberwithin{figure}{subsection} % 
%\usepackage{minted}





\begin{document}


\begin{titlepage}
    \centering
    {\Huge 电路基础讲义 \\[1cm]}
    {\Large 基于 Multisim 实现 \\[2cm]}
    {\Large 笛芦-Fumo\_B\_King \\[1cm]}
    {\large \today}
\end{titlepage}


\tableofcontents        % 生成目录
\pagenumbering{roman} 
\setcounter{page}{1}
\newpage                % 换页

\pagenumbering{arabic}
\setcounter{page}{1}
\section{电路与数学模型}

    在讲解我们的电路基础之前，我们必须要提到相关的数学模型，这并非像我们在中学阶段
    学习的物理学电路一样，我们在本专科阶段学习的电路是对一种信号与系统的具体化的工
    具，就像是我们在中学阶段学习物理学一样，物理学本身就是对真实世界的一种抽象化的
    描述。然而我们的电路是对电路这一具体的内容进行一种抽象化描述，所以在此之前我们
    是有必要对电路的抽象内容进行一点探讨的。

    我们需要知道的是，对于数学模型是有一套专门的求解方法的，读者一般来说是学习完高
    等数学和线性代数之后进修的本课程，所以对微分方程和线性方程组是有一定的了解与基
    础的。
    常用数学模型与求解方法，我们将展示电路方程的普遍形式：对于动态电路，其核心
    是常微分方程(组)；对于电阻电路，则归结为代数方程(这里可以分为线性方程组和非
    线性方程组但是研究非线性方程组是比较困难的，本课程主要以线性方程组下手)。
    我们将介绍从时域直接求解微分方程的经典法，到利用初始条件确定特解的完整过程，
    同时也会涉及数值解法在现代仿真中的意义，为后续计算机辅助分析，例如MATLAB等软件
    确定一些基础。

    在常用数学物理方法介绍中，我们将引入拉普拉斯变换与傅里叶变换这两个强力工具。通过
    将时域的微分运算转化为复频域或频域的代数运算，跳过复杂的积分算法变为简单的代数学
    算法；而复数相量法作为正弦稳态电路的研究工具，本质上是傅里叶分析在单频激励下的具
    体化，这些数学物理方法极大地降低了求解复杂度。

    在常用信号与系统方法介绍部分，我们将把电路置于信号与系统的框架中审视：电路作为系
    统，对激励信号产生响应。我们将讨论冲激响应、卷积积分、传递函数与频率响应等概念，
    阐明线性时不变电路的叠加特性与系统函数描述。这不仅是对电路工具求解的提高，更是连
    接本课程与后续“信号与系统”课程的桥梁。


    \newpage
    \subsection{电路模型介绍}

    \subsubsection{电压电流与函数关系}
    首先我们要确定两个独立的量，也可以称作为变量-电压和电流。这两个是完全独立的量，他们
    的本身是没有直接的关系的，但是在特定的元件下可以建立关系。这是因为这两个物理量的定义
    是完全独立的。电流的定义我们在中学阶段就学习过，定义为导体的横截面单位时间通过的电荷
    量。$$ I = \frac{dQ}{dt} $$
    电压的定义我们在中学阶段说的比较简单，严格的说，电压是空间中电场关于空间距离的积分。
    $$ U = \int \vec{E} \cdot d\vec{l} $$
    不过在导体中，尤其是我们最常见的线性导体中，电压和电流呈现线性变化关系。通常我们称作为
    欧姆定律。$$ U = I \cdot R $$
    其中，这里的R被称为电阻，单位为欧姆，反之，也可以写为$$ I = G \cdot U $$
    这里的G被称为电导，单位为西门子。我们不难发现电阻和电导之间的关系是倒数关系。
    $$ G = \frac{1}{R} $$
    但整体来说电压和电流是独立的量，它们之间没有直接的关系，
    但在一个系统(鉴于读者没有了解过系统的概念，可以理解为是一种特定的电路)中可以呈现出关系的
    变换。

    复杂的来说，中学阶段我们只接触过常数电压和常数电流，又或者是正弦稳态的交流电计算，但是总的
    来说都是一种常数值。读者需要注意，正弦稳态交流电并不是线性的，但是在指数代换的方法下可以使
    用复数工具进行计算。我们接下来遇到的电压或电流往往是一种函数，而不是常数值。我们可以看一种
    有趣的电压函数-冲激信号：
    $$
    \left\{
    \begin{aligned}
        u(t) = 0, t \neq 0 \\
        \int_{\mathbb{R}}  u(t) dt = 1
    \end{aligned}
    \right.
    $$

    \begin{figure}[htbp]
    \centering
    \begin{tikzpicture}[scale=1.5]
    % 坐标轴
    \draw[->] (-3,0) -- (3,0) node[right] {$t$};
    \draw[->] (0,-0.5) -- (0,2) node[above] {$\delta(t)$};
    % 冲激箭头
    \draw[->, thick, red] (0,0) -- (0,1.5) node[above right] {$\infty$};
    % 原点标记
    \fill (0,0) circle (2pt);
    \node[below left] at (0,0) {$0$};
    % 积分面积示意
    \draw[blue, dashed] (-0.3,0) -- (-0.3,0.5) -- (0.3,0.5) -- (0.3,0);
    \node[blue] at (0,0.7) {$1$};
    \end{tikzpicture}
    \caption{冲激信号函数图像}
    \end{figure}
    如果读者学习过相关内容应该知道这是一个信号与系统中的一个常见信号-冲激信号。
    这个函数显然是不可能存在的，这是一个理想化的模型，现实中冲激信号更多类似于
    量子力学中的$\delta$函数。当读者面对这种以函数出现的信号，比如电压电流时，可能会
    束手无策了，我们本课程的目的就是系统的去分析这些信号，在后续课程中，我们会介绍
    强有力的模型。
    
    希望读者应该了解电压和电流之间的关系，如果在一个电路系统中，我们知道电压函数与电流函数
    我们可以推导出这个电路系统所具备的性质，这是我们后续课程中的传递阻抗与传递导纳的内容了。

    \subsubsection{电阻与线性函数}
    在中学阶段我们最常见的就是纯电阻电路了，中学老师说的很多定理和推到的前提都是这个
    电路是一个纯电阻电路。那至于为什么是纯电阻电路，就是因为纯电阻电路是一个完全的线性
    系统。这里面所有的电压和电流全部呈现线性变化。比如我们发现在电阻完全恒定的情况下
    随着电流大小的变化，电压也会随着线性的变化。让我们对比一下正比函数与欧姆定律：
    $$
    U = R \cdot I   \\
    y = k \cdot x
    $$
    我们发现这里面常数R完完全全就是一个正比函数。因为电阻往往是常数所以它是一个完全的
    线性函数。不过在中学阶段读者们会有一个疑惑，比如电阻如果变化那还怎么求解？甚至这是
    一个"完完全全的悖论"，为什么这么说？因为：电压越大电流越大，电流越大灯泡越亮，灯泡
    越亮灯泡越热，灯泡越热电阻越大，电阻越大电流越小，于是得出一个结论，电流越大电流越
    小。但是想想也不可能是这样。这里的问题出自哪里呢？
    $$
        i = \frac{u}{R} \\
        R = f(i)
    $$
    我们发现这里面的电阻和电流之间还有一层关系，这个关系目前是不得而知的，若读者有兴趣
    可以查询相关资料了解电流、温度和电阻的关系。这里的电阻都是变化的，所以这就不再是一
    个简单的线性函数了。所以我们要求的纯电阻电路是线性的要求就是电阻的值不能改变。

    让我们研究一下纯电阻的系统：
    \begin{figure}[htbp]
    \centering
    \begin{circuitikz}

    \draw (0,0) to [short, o-] (2,0)
                to[R, l = $R$] (4,0)
                to [short, -o] (6,0);
    \node[below] at (0,0) {$+$};
    \node[below] at (6,0) {$-$};
    \node[below] at (3,-1) {$U$};
    
    \draw[->] (2, 1) -- (4, 1);
    \node[above] at (3,1) {$I$};

    \end{circuitikz}
    \caption{电阻电路}
    \end{figure}




    \newpage
    \subsection{常用数学模型与求解方法}

    \newpage
    \subsection{常用数学物理方法介绍}

    \newpage
    \subsection{常用信号与系统方法介绍}

    \newpage
    \subsection{习题与详解}


\end{document}